% ════════════════════════════════════════════════════════════════════════════════════════
\section{Parte 1: Introducción}

\begin{frame}[fragile]{¿Qué es Python?}
  \begin{columns}[T]
    \begin{column}{0.52\linewidth}
      \begin{itemize}\setlength\itemsep{4pt}
        \item Un \alert{lenguaje de programación de propósito general}: gratuito, legible y con una enorme comunidad.
        \item Describes \emph{qué hacer} paso a paso y el ordenador lo hace, con 10 muestras o con 10 millones.
        \item En ciencia es el \alert{pegamento} entre los datos y el análisis. Paquetes que usamos hoy:
      \end{itemize}
      \small
      \begin{tabular}{@{}ll@{}}
        \code{numpy} & matrices y álgebra lineal\\
        \code{pandas} & tablas\\
        \code{scipy}, \code{statsmodels} & estadística\\
        \code{matplotlib} & gráficos\\
        \code{scikit-learn} & aprendizaje automático\\
        \code{anndata} & contenedores de datos ómicos\\
      \end{tabular}
    \end{column}
    \begin{column}{0.46\linewidth}
\begin{lstlisting}
genes  = ["STAT1", "CXCL9", "ACTB"]
counts = {"STAT1": 73853, "ACTB": 36500}
total  = sum(counts.values())
print(f"{total:,} UMI en total")
\end{lstlisting}
      \vspace{0.15cm}
      \begin{block}{Hoy no necesitas ser programador}
        Te damos el código. Tu trabajo: \textbf{ejecutarlo, leer el resultado y entender qué significan los números}.
      \end{block}
    \end{column}
  \end{columns}
\end{frame}

\begin{frame}{¿Qué es un notebook de Jupyter?}
  \begin{columns}[c]
    \begin{column}{0.56\linewidth}
      \centering
      \begin{tikzpicture}[font=\scriptsize, every node/.style={align=left}]
        \node[draw=greyink!60,fill=white,minimum width=6.6cm,minimum height=3.9cm,anchor=north west] (page) at (0,0) {};
        % markdown cell
        \node[draw=none,fill=softgrey,minimum width=6.2cm,minimum height=0.95cm,anchor=north west] (md) at (0.2,-0.2) {};
        \node[anchor=north west,font=\small\bfseries] at (0.35,-0.28) {Parte 2: normalización};
        \node[anchor=north west,font=\scriptsize,text=greyink] at (0.35,-0.68) {Los conteos crudos no son comparables\ldots};
        % code cell
        \node[draw=bluemid,fill=bluemid!6,minimum width=6.2cm,minimum height=0.95cm,anchor=north west] (code) at (0.2,-1.3) {};
        \node[anchor=north west,font=\ttfamily\scriptsize] at (0.35,-1.42) {adata = ad.read\_h5ad("data.h5ad")};
        \node[anchor=north west,font=\ttfamily\scriptsize] at (0.35,-1.78) {print(adata.shape)};
        \node[circle,fill=bluemid,text=white,inner sep=2pt,font=\tiny\bfseries] at (0.02,-1.8) {$\blacktriangleright$};
        % output
        \node[anchor=north west,font=\ttfamily\scriptsize,text=ucoblue] at (0.35,-2.45) {(4224, 18120)};
        \node[anchor=north west,font=\scriptsize,text=greyink] at (3.3,-2.45) {$\leftarrow$ el resultado};
        % labels
        \node[anchor=west,text=accentred,font=\scriptsize\bfseries] at (6.9,-0.65) {texto};
        \node[anchor=west,text=accentred,font=\scriptsize\bfseries] at (6.9,-1.75) {código};
        \node[anchor=west,text=accentred,font=\scriptsize\bfseries] at (6.9,-2.55) {resultado};
        \draw[accentred,-{Stealth}] (6.85,-0.65) -- (6.45,-0.65);
        \draw[accentred,-{Stealth}] (6.85,-1.75) -- (6.45,-1.75);
        \draw[accentred,-{Stealth}] (6.85,-2.55) -- (6.45,-2.55);
      \end{tikzpicture}
    \end{column}
    \begin{column}{0.42\linewidth}\small
      \begin{itemize}\setlength\itemsep{3pt}
        \item Un documento que mezcla \alert{texto, código y resultados}.
        \item Dos tipos de celdas: de \textbf{texto} (Markdown) y de \textbf{código} (Python).
        \item Ejecutar una celda: \alert{Shift\,+\,Enter}. El resultado aparece debajo.
        \item Un \textbf{kernel} (el proceso de Python que hay detrás) \alert{recuerda} las variables: ejecuta las celdas \textbf{en orden}.
        \item Reiniciar el kernel $\Rightarrow$ se borra la memoria.
      \end{itemize}
    \end{column}
  \end{columns}
\end{frame}

\begin{frame}{¿Cómo lo vamos a ejecutar? Google Colab}
  \begin{columns}[c]
    \begin{column}{0.60\linewidth}
      \centering
      \begin{tikzpicture}[font=\scriptsize, every node/.style={align=left}]
        \draw[rounded corners=3pt,draw=greyink!70,fill=white] (0,0) rectangle (7.6,-4.2);
        \fill[greyink!12] (0,0) rectangle (7.6,-0.55);
        \node[anchor=west,font=\scriptsize\bfseries] at (0.15,-0.28) {Practica\_v3\_es.ipynb};
        \node[anchor=east,draw=bluemid,fill=bluemid!10,rounded corners=2pt,inner sep=2pt,font=\tiny] at (7.45,-0.28) {Conectar};
        \fill[greyink!5] (0,-0.55) rectangle (7.6,-0.95);
        \node[anchor=west,font=\tiny] at (0.15,-0.75) {Archivo\ \ Editar\ \ Ver\ \ Insertar\ \ \textbf{Entorno}\ \ Herram.\ \ Ayuda};
        \node[anchor=west,font=\tiny,draw=greyink!50,inner sep=1.5pt] at (4.7,-0.75) {+ Código};
        \node[anchor=west,font=\tiny,draw=greyink!50,inner sep=1.5pt] at (5.75,-0.75) {+ Texto};
        \fill[greyink!8] (0,-0.95) rectangle (1.45,-4.2);
        \node[anchor=north west,font=\tiny] at (0.1,-1.05) {$\equiv$ Índice\\ Parte 1\\ Parte 2\\ \quad 2.1 \ldots\\ Parte 3\\ Parte 4};
        \node[draw=bluemid,fill=bluemid!6,minimum width=5.6cm,minimum height=1.15cm,anchor=north west] at (1.75,-1.15) {};
        \node[anchor=north west,font=\ttfamily\tiny] at (2.2,-1.25) {adata = ad.read\_h5ad(DATA\_PATH)};
        \node[circle,fill=bluemid,text=white,inner sep=2.5pt,font=\tiny] at (1.95,-1.7) {$\blacktriangleright$};
        \node[anchor=north west,font=\ttfamily\tiny,text=ucoblue] at (2.2,-1.75) {\checkmark\ Cargado en memoria};
        \node[draw=none,fill=softgrey,minimum width=5.6cm,minimum height=0.8cm,anchor=north west] at (1.75,-2.6) {};
        \node[anchor=north west,font=\scriptsize\bfseries] at (1.9,-2.68) {2.2 Un número cada vez};
        \node[anchor=north west,font=\tiny,text=greyink] at (1.9,-3.02) {las celdas de texto explican, las de código calculan};
        % numbered callouts
        \foreach \n/\x/\y/\tx/\ty in {1/7.2/0.28/7.2/-0.12, 2/1.55/-1.0/1.55/-1.3, 3/3.9/0.28/3.9/-0.55, 4/-0.3/-1.6/0.05/-1.6} {
          \node[circle,fill=accentred,text=white,inner sep=1.5pt,font=\tiny\bfseries] (c\n) at (\x,\y) {\n};
          \draw[accentred,thick] (c\n) -- (\tx,\ty);
        }
      \end{tikzpicture}
    \end{column}
    \begin{column}{0.38\linewidth}\small
      \begin{itemize}\setlength\itemsep{3pt}
        \item Jupyter en tu \alert{navegador}: nada que instalar. Gratis, con una cuenta de Google.
        \item Tu código se ejecuta en un \alert{ordenador temporal} en la nube de Google.
      \end{itemize}
      \vspace{0.1cm}
      \begin{enumerate}\setlength\itemsep{2pt}\footnotesize
        \item[\textcolor{accentred}{\textbf{1}}] \textbf{Conectar} a una sesión
        \item[\textcolor{accentred}{\textbf{2}}] $\blacktriangleright$ \textbf{ejecutar} una celda
        \item[\textcolor{accentred}{\textbf{3}}] menú \textbf{Entorno de ejecución}: ejecutar todo, reiniciar
        \item[\textcolor{accentred}{\textbf{4}}] \textbf{Índice}: saltar entre partes
      \end{enumerate}
    \end{column}
  \end{columns}
\end{frame}

\begin{frame}{Kit de supervivencia de Colab}
  \begin{columns}[T]
    \begin{column}{0.52\linewidth}
      \small
      \begin{tabular}{@{}>{\ttfamily}p{0.38\linewidth}p{0.58\linewidth}@{}}
        \toprule
        \textrm{\textbf{Atajo}} & \textbf{Qué hace}\\\midrule
        Shift+Enter & ejecuta y pasa a la siguiente\\
        Ctrl+Enter & ejecuta y se queda\\
        Ctrl+M B & inserta una celda debajo\\
        Ctrl+M D & borra la celda\\
        \bottomrule
      \end{tabular}
      \vspace{0.3cm}

      \textbf{Si algo se rompe}
      \begin{itemize}\footnotesize\setlength\itemsep{2pt}
        \item \textbf{Entorno de ejecución $\rightarrow$ Reiniciar sesión} y vuelve a ejecutar desde el principio.
        \item ¿Te has quedado atrás en la Parte 2? Usa la \alert{celda de rescate} al inicio de la Parte 3.
      \end{itemize}
    \end{column}
    \begin{column}{0.44\linewidth}
      \begin{block}{Conviene saber}
        \footnotesize
        \begin{itemize}\setlength\itemsep{3pt}
          \item Las sesiones son \alert{temporales}: archivos y variables desaparecen al terminar. Descargamos los datos al principio (alrededor de un minuto).
          \item Una sesión inactiva se desconecta: reconecta y vuelve a ejecutar.
          \item Las celdas cuyo título termina en \emph{``ejecutar una vez, no hace falta leer''} están plegadas a propósito: son código auxiliar para los gráficos.
          \item Para conservar tus cambios: \textbf{Archivo $\rightarrow$ Guardar una copia en Drive}.
        \end{itemize}
      \end{block}
    \end{column}
  \end{columns}
\end{frame}

% ────────────────────────────────────────────────────────────────────────────────────────
\begin{frame}{¿Qué es la transcriptómica mediante RNA-seq?}
  \centering
  \begin{tikzpicture}[font=\scriptsize, every node/.style={align=center},
      stp/.style={draw=ucoblue,fill=ucoblue!8,rounded corners,minimum height=1.35cm,minimum width=2.05cm,inner sep=2pt}]
    \node[stp] (a) {tejido o\\ células};
    \node[stp,right=0.35cm of a] (b) {\textbf{extracción}\\ de RNA};
    \node[stp,right=0.35cm of b] (c) {\textbf{librería}\\ fragmentos +\\ adaptadores};
    \node[stp,right=0.35cm of c] (d) {\textbf{secuenciación}\\ millones de\\ lecturas};
    \node[stp,right=0.35cm of d] (e) {\textbf{alineamiento}\\ lecturas $\to$ genes};
    \node[stp,right=0.35cm of e,fill=accentred!10,draw=accentred] (f) {\textbf{matriz de}\\ \textbf{conteos}};
    \foreach \x/\y in {a/b,b/c,c/d,d/e,e/f} \draw[-{Stealth},thick] (\x) -- (\y);
  \end{tikzpicture}
  \vspace{0.35cm}

  \begin{columns}[T]
    \begin{column}{0.5\linewidth}\small
      \begin{itemize}\setlength\itemsep{3pt}
        \item El \alert{transcriptoma} es el conjunto de todas las moléculas de RNA de una muestra: qué genes están \emph{encendidos} y cuánto.
        \item El \textbf{RNA-seq bulk} mide \alert{un perfil promedio por muestra}: todas las células están mezcladas.
        \item Un \textbf{conteo} es el número de moléculas (o lecturas) asignadas a un gen.
      \end{itemize}
    \end{column}
    \begin{column}{0.46\linewidth}
      \centering\footnotesize
      \begin{tabular}{lrrr}
        \toprule
        gen & muestra 1 & muestra 2 & muestra 3\\\midrule
        \gene{STAT1} & 5\,310 & 12\,874 & 4\,990\\
        \gene{CXCL9} & 12 & 1\,950 & 9\\
        \gene{ACTB} & 36\,500 & 41\,200 & 33\,800\\
        $\vdots$ & $\vdots$ & $\vdots$ & $\vdots$\\
        \bottomrule
      \end{tabular}\\[2pt]
      \scriptsize una matriz de conteos: \alert{genes $\times$ muestras} (números ilustrativos)
    \end{column}
  \end{columns}
\end{frame}

\begin{frame}{RNA-seq bulk y de célula única: los dos flujos de trabajo}
  \begin{columns}[c]
    \begin{column}{0.3\linewidth}
      \centering\includegraphics[width=\linewidth,height=0.8\textheight,keepaspectratio]{tzec_fig1_bulk.png}
    \end{column}
    \begin{column}{0.3\linewidth}
      \centering\includegraphics[width=\linewidth,height=0.8\textheight,keepaspectratio]{tzec_fig2_sc.png}
    \end{column}
    \begin{column}{0.36\linewidth}\small
      \begin{itemize}\setlength\itemsep{4pt}
        \item La misma lógica con distinta unidad: una \alert{muestra} (bulk) o una \alert{célula} (célula única).
        \item Ambos terminan en una \alert{matriz de conteos}, que luego se normaliza y se analiza (expresión diferencial, enriquecimiento\ldots).
        \item La célula única añade \alert{códigos de barras} para asignar cada lectura a su célula, y control de calidad por célula.
      \end{itemize}
      {\tiny Figuras 1 y 2 de Tzec-Interi\'an \emph{et al.}\ 2025, \emph{Quant.\ Biol.}\ (CC BY)}
    \end{column}
  \end{columns}
\end{frame}

\begin{frame}{Transcriptómica de célula única: un perfil por célula}
  \begin{columns}[c]
    \begin{column}{0.5\linewidth}
      \centering
      \begin{tikzpicture}[font=\scriptsize, every node/.style={align=center}]
        % bulk
        \node[font=\scriptsize\bfseries] at (1.0,2.6) {bulk};
        \foreach \c in {0,...,3} { \fill[ucoblue!\the\numexpr 35+\c*12\relax,draw=white] (\c*0.5,2.1) rectangle ++(0.45,-0.35); }
        \node[font=\tiny,text=greyink] at (1.0,1.45) {1 muestra = 1 fila: promedio de todas las células};
        % single cell
        \node[font=\scriptsize\bfseries] at (1.0,0.85) {célula única};
        \foreach \r in {0,...,5} {
          \foreach \c in {0,...,3} {
            \pgfmathsetmacro{\v}{ifthenelse(rnd>0.72,rnd*65+35,0)}
            \fill[bluemid!\v!white,draw=greyink!25] (\c*0.5,0.4-\r*0.28) rectangle ++(0.45,-0.25);
          }
        }
        \node[font=\tiny,text=greyink] at (1.0,-1.55) {1 célula = 1 fila: miles de filas, casi todo ceros};
      \end{tikzpicture}
    \end{column}
    \begin{column}{0.48\linewidth}\small
      \begin{tabular}{@{}>{\raggedright\arraybackslash}p{0.2\linewidth}>{\raggedright\arraybackslash}p{0.25\linewidth}>{\raggedright\arraybackslash}p{0.45\linewidth}@{}}
        \toprule
        & \textbf{bulk} & \textbf{célula única}\\\midrule
        unidad & muestra & célula\\
        filas & decenas & $10^4$--$10^7$\\
        ceros & pocos & casi todas las entradas\\
        revela & el promedio & tipos celulares y células raras\\
        \bottomrule
      \end{tabular}
      \vspace{0.25cm}
      \begin{itemize}\footnotesize\setlength\itemsep{3pt}
        \item Cada célula recibe un \alert{código de barras}, de modo que las lecturas se asignan a su célula.
        \item Los tipos celulares se identifican después por clustering y \emph{genes marcadores} (p.\,ej.\ \gene{LYZ} en monocitos, \gene{CD3E} en linfocitos T).
      \end{itemize}
    \end{column}
  \end{columns}
\end{frame}

\begin{frame}{Célula única por indexado combinatorio (split-pool)}
  \centering
  \begin{tikzpicture}[scale=0.92,transform shape,font=\scriptsize, every node/.style={align=center}]
    % plates: 3 small grids
    \foreach \p/\x/\col in {1/0.3/ucoblue, 2/3.5/orangeslot, 3/6.7/greenslot} {
      \foreach \r in {0,...,2} \foreach \c in {0,...,3} {
        \pgfmathsetmacro{\sh}{20+mod((\r*4+\c)*7+\p*11,70)}
        \fill[\col!\sh!white] (\x+\c*0.3,-\r*0.3) rectangle ++(0.26,-0.26);
      }
      \node[font=\scriptsize\bfseries,text=\col] at (\x+0.6,0.3) {ronda \p};
    }
    \node[draw=greyink!50,rounded corners,minimum width=1.9cm,minimum height=0.8cm,font=\scriptsize] (cells) at (-1.5,-0.45) {células fijadas\\ \tiny (cada célula es su ``tubo'')};
    \draw[-{Stealth},thick] (cells) -- (0.1,-0.45);
    \draw[-{Stealth},thick] (1.65,-0.45) -- node[above,font=\tiny] {mezclar y repartir} (3.3,-0.45);
    \draw[-{Stealth},thick] (4.85,-0.45) -- node[above,font=\tiny] {mezclar y repartir} (6.5,-0.45);
    \node[draw=accentred,fill=accentred!8,rounded corners,minimum width=2.5cm,minimum height=0.95cm] (res) at (10.2,-0.45) {cada célula: un código\\ único \textbf{B1\,--\,B2\,--\,B3}};
    \draw[-{Stealth},thick] (8.05,-0.45) -- (res);
  \end{tikzpicture}
  \vspace{0.3cm}

  \begin{columns}[T]
    \begin{column}{0.54\linewidth}\small
      \begin{itemize}\setlength\itemsep{3pt}
        \item En cada ronda las células se \alert{reparten en pocillos al azar}; cada pocillo añade su propio código de barras al RNA de las células que contiene.
        \item Tras tres rondas, la probabilidad de que dos células compartan los mismos tres códigos es minúscula: con 96 códigos por ronda, $96^3\approx 880\,000$ combinaciones (ilustrativo).
        \item \textbf{Sin gotas ni aislamiento de células}: millones de células en un solo experimento.
      \end{itemize}
    \end{column}
    \begin{column}{0.42\linewidth}
      \begin{block}{En nuestros datos}
        \footnotesize
        El \alert{código de la primera ronda} es el pocillo de la placa de estimulación, es decir, la \alert{condición}: donante $\times$ citoquina. Es la columna \code{bc1\_well} del archivo crudo.\\[2pt]
        El control PBS se repartió en \alert{6 pocillos}.
      \end{block}
    \end{column}
  \end{columns}
\end{frame}

\begin{frame}{¿Por qué hay que hacer pseudobulk primero?}
  \centering
  \begin{tikzpicture}[font=\scriptsize]
    \foreach \r in {0,...,5} {
      \foreach \c in {0,...,7} {
        \pgfmathsetmacro{\v}{ifthenelse(rnd>0.5,rnd*65+35,0)}
        \fill[bluemid!\v!white,draw=greyink!25] (\c*0.28,-\r*0.28) rectangle ++(0.26,-0.26);
      }
    }
    \node[align=center] at (1.1,0.6) {células de un donante, una citoquina, un tipo celular\\ \textbf{células $\times$ genes} (conteos UMI)};
    \draw[decorate,decoration={brace,amplitude=4pt,mirror}] (-0.1,0) -- (-0.1,-1.65) node[midway,left=5pt] {células};
    \draw[-{Stealth},very thick,accentred] (2.5,-0.8) -- node[above] {$\sum$ sobre células} (4.5,-0.8);
    \foreach \c in {0,...,7} {
      \pgfmathsetmacro{\v}{rnd*45+50}
      \fill[ucoblue!\v] (4.8+\c*0.28,-0.55) rectangle ++(0.26,-0.26);
    }
    \node[align=center] at (5.9,0.2) {\textbf{una fila}\\ un perfil ``tipo bulk''};
    \draw[thick] (8.6,0.2) rectangle (11.6,-1.9);
    \foreach \r in {1,...,6} \draw[gray!50] (8.6,0.2-0.3*\r) -- (11.6,0.2-0.3*\r);
    \fill[ucoblue!60] (8.6,-0.7) rectangle (11.6,-1.0);
    \node[right] at (11.7,-0.85) {1 muestra};
    \node at (10.1,-2.3) {\textbf{X}: \numnsamples\ muestras $\times$ \numngenesraw\ genes};
  \end{tikzpicture}
  \vspace{0cm}

  \begin{itemize}\small\setlength\itemsep{2pt}
    \item \alert{La unidad de replicación es el donante, no la célula}: millones de células proceden de solo 12 personas, y las células de una persona no son independientes. Analizar las células como si fueran réplicas produce muchos falsos positivos (\emph{pseudorreplicación}).
    \item Los conteos de célula única son dispersos y ruidosos; sumados sobre miles de células se comportan como \alert{RNA-seq bulk}, así que se aplica una estadística estándar y transparente.
    \item Hace los datos pequeños: 9.7 millones de células $\to$ \numnsamples\ muestras.
  \end{itemize}
\end{frame}

% ────────────────────────────────────────────────────────────────────────────────────────
\begin{frame}{El conjunto de datos: el Diccionario Humano de Citoquinas}
  \begin{columns}[c]
    \begin{column}{0.52\linewidth}
      \centering
      \includegraphics[width=\linewidth,height=0.74\textheight,keepaspectratio]{paper_fig1ab.png}\\[-1pt]
      {\tiny Oesinghaus \emph{et al.}\ 2025, Fig.\ 1a--b (CC BY-NC-ND 4.0)}
    \end{column}
    \begin{column}{0.44\linewidth}\small
      \begin{itemize}\setlength\itemsep{4pt}
        \item \alert{12 donantes sanos} (6 mujeres, 6 hombres; 34--75 años)
        \item \alert{90 citoquinas} + un control PBS, in vitro durante \alert{24 h}
        \item $\sim$\alert{9.7 millones} de PBMC tras el control de calidad
        \item 16 tipos celulares, anotados con genes marcadores (panel b)
      \end{itemize}
      \vspace{0.2cm}
      \footnotesize
      \textbf{Pregunta del artículo:} ¿qué hace cada citoquina a cada tipo celular inmune y cuánto difieren las personas?
    \end{column}
  \end{columns}
\end{frame}

\begin{frame}{Cómo se generaron los datos}
  \centering
  \begin{tikzpicture}[font=\scriptsize, every node/.style={align=center},
      stp/.style={draw=ucoblue,fill=ucoblue!8,rounded corners,minimum height=1.45cm,minimum width=1.85cm,inner sep=2pt},
      numdot/.style={circle,fill=accentred,text=white,inner sep=1.5pt,font=\tiny\bfseries}]
    \node[stp] (a) {PBMC de\\ 12 donantes};
    \node[stp,right=0.3cm of a] (b) {placa de 96\\ 10$^6$ células/pocillo\\ citoquina o PBS\\ \textbf{24 h}};
    \node[stp,right=0.3cm of b] (c) {fijar las\\ células};
    \node[stp,right=0.3cm of c] (d) {\textbf{3 rondas}\\ de barcoding\\ split-pool};
    \node[stp,right=0.3cm of d] (e) {secuenciación\\ $\sim$31\,000\\ lecturas/célula};
    \node[stp,right=0.3cm of e,fill=accentred!10,draw=accentred] (f) {conteos\\ célula $\times$ gen};
    \foreach \x/\y in {a/b,b/c,c/d,d/e,e/f} \draw[-{Stealth},thick] (\x) -- (\y);
    \node[stp,below=0.75cm of d,fill=white] (g) {control de calidad\\ eliminar dobletes};
    \node[stp,left=0.3cm of g,fill=white] (h) {clustering y\\ anotación:\\ 16 tipos celulares};
    \node[stp,right=0.3cm of g,fill=white,draw=accentred] (i) {\textbf{pseudobulk}\\ donante $\times$ citoquina\\ $\times$ tipo celular};
    \draw[-{Stealth},thick] (f.south) |- ($(g.north)+(0.5,0.35)$) -- ($(g.north)+(0.5,0)$);
    \draw[-{Stealth},thick] (g) -- (i); \draw[-{Stealth},thick] (h) -- (g);
  \end{tikzpicture}
  \vspace{0.25cm}

  \footnotesize
  \begin{itemize}\setlength\itemsep{2pt}
    \item Barcoding con el protocolo GigaLab de Parse Biosciences; sublibrerías de $\sim$100\,000 células; secuenciación en un Ultima UG-100.
    \item Se conservaron células con 400--7\,000 genes detectados y $<15\,\%$ de lecturas mitocondriales.
  \end{itemize}
\end{frame}

\begin{frame}{¿Qué contiene el subconjunto con el que trabajamos?}
  \begin{columns}[T]
    \begin{column}{0.56\linewidth}
      \small
      \begin{tabular}{@{}>{\raggedright\arraybackslash}p{0.2\linewidth}>{\raggedright\arraybackslash}p{0.22\linewidth}>{\raggedright\arraybackslash}p{0.52\linewidth}@{}}
        \toprule
        & \textbf{artículo} & \textbf{nuestra práctica}\\\midrule
        donantes & 12 & 12\\[1pt]
        citoquinas & 90 + PBS & 87 + PBS\\[1pt]
        tipos celulares & 16 & 4: T CD4, T CD8, NK CD56\textsuperscript{low}, monocitos CD14\\[1pt]
        unidad & células únicas & muestras \alert{pseudobulk}\\[1pt]
        genes & todos & \numngenesraw\ codificantes de proteínas ($\to$ \numngenes\ tras filtrar)\\
        \bottomrule
      \end{tabular}
    \end{column}
    \begin{column}{0.4\linewidth}
      \begin{block}{El diseño}
        \centering
        $4$ tipos celulares\\ $\times$ $88$ condiciones\\ \footnotesize(87 citoquinas + PBS)\\ \normalsize $\times$ $12$ donantes\\[3pt]
        $=$ \alert{\numnsamples\ muestras}
      \end{block}
      \footnotesize
      El control PBS se midió en 6 pocillos; el archivo guarda cada pocillo \emph{y} su total. Nos quedamos con el total (extra B1 del notebook).
    \end{column}
  \end{columns}
\end{frame}

\begin{frame}{AnnData: una estructura de datos de Python para datos ómicos}
  \begin{columns}[c]
    \begin{column}{0.38\linewidth}
      \centering
      \includegraphics[width=\linewidth,height=0.76\textheight,keepaspectratio]{anndata_schema.png}\\[-1pt]
      {\tiny Esquema de AnnData (documentación de anndata)}
    \end{column}
    \begin{column}{0.58\linewidth}\small
      \begin{itemize}\setlength\itemsep{3pt}
        \item \textbf{Un objeto con varias tablas que siempre permanecen alineadas.}
        \item \textcolor{ucoblue}{\textbf{\code{X}}}: la matriz principal, $n_{\text{obs}}\times n_{\text{vars}}$
        \item \textcolor{yellowslot}{\textbf{\code{obs}}}: una fila por \alert{observación} (muestra o célula): metadatos
        \item \textcolor{bluemid}{\textbf{\code{var}}}: una fila por \alert{variable} (gen): metadatos
        \item \textcolor{greenslot}{\textbf{\code{layers}}}: más matrices, con la misma forma que \code{X}
        \item \code{obsm}, \code{varm}, \code{obsp}, \code{varp}, \code{uns}: embeddings (PCA), grafos, resultados de formato libre
      \end{itemize}
      \vspace{0.15cm}
      \begin{block}{Por qué importa}
        \footnotesize Seleccionar un subconjunto mueve \emph{todas} las tablas a la vez: ningún gen ni muestra puede desalinearse de sus metadatos. Se guarda como archivo \code{.h5ad}.
      \end{block}
    \end{column}
  \end{columns}
\end{frame}

\begin{frame}[fragile]{AnnData en nuestra práctica}
  \begin{columns}[T]
    \begin{column}{0.52\linewidth}
      \small
      \begin{tabular}{@{}>{\ttfamily}p{0.16\linewidth}>{\raggedright\arraybackslash}p{0.47\linewidth}>{\raggedright\arraybackslash}p{0.3\linewidth}@{}}
        \toprule
        \textrm{\textbf{slot}} & \textbf{nuestros datos} & \textbf{forma}\\\midrule
        X & conteos crudos, luego $\log_2$(CPM+1) & \numnsamples\,$\times$\,\numngenesraw\\[2pt]
        obs & donor, cytokine, treatment, cell\_type, número de células & \numnsamples\ filas\\[2pt]
        var & nombres de genes, conteos totales & \numngenesraw\ filas\\[2pt]
        layers & \code{counts}, \code{cpm} & como \code{X}\\[2pt]
        obsm & puntuaciones del PCA (Parte 3) & \numnsamples\,$\times$\,10\\
        \bottomrule
      \end{tabular}
    \end{column}
    \begin{column}{0.46\linewidth}
\begin{lstlisting}
# primero filas (muestras), luego genes
adata[["Donor1|CD14_Mono|IFN-gamma"],
      ["STAT1"]]

mask = adata.obs["cell_type"] == "CD14_Mono"
adata[mask, ["IDO1", "CXCL10"]]
\end{lstlisting}
      \footnotesize La primera selección es de $1\times1$; la segunda, de $1\,056\times2$ (todas las muestras de monocitos).
    \end{column}
  \end{columns}
  \vspace{0.25cm}
  \begin{block}{Manos a la obra}
    \small Volvemos al notebook: \textbf{preparación de la Parte 1} (instalar, cargar los datos) y después la \textbf{Parte 2}.
    \timebox{unos 25 min para la Parte 1}
  \end{block}
\end{frame}
